\chapter{Аналитический обзор}
\label{ch:chap1}

\section{Визуально-языковые модели действий}
\label{sec:vla}

    Визуально-языковая модель действий (Vision-Language-Action model, VLA) -- это модель, использующая предобученный backbone, обученный на крупномасштабных визуально-языковых данных из Интернета, и впоследствии дообученная на генерацию управляющих команд робота \cite{reuss2025blog}. Управляющие команды могут представлять собой углы сочленений, положение концевого эффектора, состояние захвата и другие управляющие воздействия. В отличие от мультимодальных политик, использующих раздельно предобученные текстовый и визуальный энкодеры, VLA-модели опираются на совместное визуально-языковое предобучение, что теоретически обеспечивает более глубокое понимание языковых инструкций и лучшее обобщение.

    По архитектуре генерации действий современные VLA-модели подразделяются на три основных класса.

    Авторегрессионные модели представляют действия в виде последовательности дискретных токенов и генерируют их поочерёдно, аналогично генерации текста в языковых моделях. К этому классу относится RT-2 \cite{brohan2023rt2} -- модель с 55 миллиардами параметров, дообученная на робототехнических траекториях поверх предобученной визуально-языковой модели PaLI-X. На 6000 испытаниях RT-2 продемонстрировала двукратное превосходство над RT-1 при работе с новыми объектами и средами, а также emergent-способности к chain-of-thought рассуждениям. OpenVLA \cite{kim2024openvla} -- открытая 7B модель на основе Llama~2 с объединённым визуальным энкодером DINOv2 и SigLIP. Обучена на 970~000 реальных робототехнических траекторий из набора Open X-Embodiment, охватывающих более 70 доменов задач. Превосходит RT-2-X на 16,5\% по абсолютному success rate на 29 задачах, при этом имея в 8 раз меньше параметров. Поддерживает эффективное дообучение через LoRA на потребительских GPU.

    Диффузионные и flow-matching модели генерируют непрерывные траектории действий через итеративный процесс денойзинга. Модель $\pi_0$ \cite{black2024pi0} от Physical Intelligence использует PaliGemma (3B) в качестве визуально-языкового backbone и flow matching для генерации действий. Предобучена на данных с 7 конфигураций роботов и 68 задач, код и веса доступны через репозиторий openpi; на LIBERO достигает среднего SR 94,2\% (Spatial 96,8\%, Object 98,8\%, Goal 95,8\%, Long 85,2\%). Вариант $\pi_0$-FAST \cite{black2024pi0} применяет частотное токенизирование действий (FAST) для ускорения генерации и показывает 85,5\% в среднем. FLOWER \cite{reuss2025flower} -- компактная модель с 950M параметров, достигающая state-of-the-art на CALVIN ABC (4,53) при обучении всего за 200~H100-часов. Использует intermediate-modality fusion с сокращением LLM-слоёв на 50\% и перераспределением ёмкости в диффузионную голову.

    Дискретно-диффузионные модели -- новейший класс, представленный на ICLR 2026. В отличие от авторегрессионных, они генерируют последовательности действий параллельно через дискретную диффузию, что позволяет одновременно создавать и рассуждения, и действия. К этому классу относятся Discrete Diffusion VLA, dVLA и DiVA, показывающие результаты 95--98\% на LIBERO.

\subsection{Дополнительные открытые VLA-модели}
\label{subsec:more_vla}

    Помимо базовых моделей, для экспериментального исследования (глава~\ref{ch:chap3}) привлечён ряд более новых открытых VLA, демонстрирующих конкурентные результаты на LIBERO.

    \textbf{UniVLA} \cite{univla2025} -- унифицированная мультимодальная модель, обучающая зрение, язык и действие в едином авторегрессионном пространстве дискретных токенов с включением видеоданных в обучение мира. Достигает среднего SR 95,5\% на LIBERO, причём ключевой вклад -- скачок на длинногоризонтном наборе Long с прежних 69,0\% до 94,0\%.

    \textbf{CogACT} \cite{cogact2024} -- модель, разделяющая когнитивный (VLM-рассуждение) и моторный (диффузионный action-эксперт) модули. На LIBERO показывает средний SR 93,6\% (Spatial 97,2\%, Object 98,0\%, Goal 90,2\%, Long 88,8\%), демонстрируя устойчивость на длинных задачах за счёт специализированной action-головы.

    \textbf{SpatialVLA} \cite{spatialvla2025} -- VLA с явным кодированием пространственных представлений (Ego3D position encoding, adaptive action grids), backbone 4B. Средний SR на LIBERO -- 78,1\% (Spatial 88,2\%, Object 89,9\%, Goal 78,6\%, Long 55,5\%); модель сильна в пространственном обобщении, но уступает на длинногоризонтных композиционных задачах.

    \textbf{Octo} \cite{octo2024} и \textbf{OpenVLA} \cite{kim2024openvla} в исходных (не-OFT) версиях служат нижними границами: 75,1\% и 76,5\% в среднем на LIBERO соответственно, что подчёркивает значимость рецепта дообучения OFT (parallel decoding + action chunking + L1-регрессия), поднимающего OpenVLA до 97,1\%.

\subsection{Фронтирные модели}
\label{subsec:frontier}

    Отдельного рассмотрения заслуживают закрытые фронтирные модели, определяющие верхнюю границу возможностей VLA. PaLM-E \cite{driess2023palme} -- мультимодальная языковая модель с 562 миллиардами параметров, инъецирующая непрерывные сенсорные данные в пространство эмбеддингов языковой модели. Достигла state-of-the-art на задачах визуального вопросно-ответного диалога (OK-VQA) при сохранении робототехнических способностей. Gemini Robotics 1.5 \cite{deepmind2025gemini}, представленная DeepMind в сентябре 2025 года, ввела механизм motion transfer для zero-shot переноса навыков между разными платформами роботов и <<embodied thinking>> -- чередование внутренних рассуждений на естественном языке с генерацией действий.

    Между открытыми и фронтирными моделями существует разрыв, который не виден из симуляционных бенчмарков. Открытые модели, такие как FLOWER и OpenVLA, конкурентоспособны на LIBERO и CALVIN, однако значительно уступают Gemini Robotics и $\pi_{0.5}$ по способности к zero-shot обобщению на произвольные объекты и инструкции в реальном мире \cite{reuss2025blog}. Среди причин этого разрыва -- ограниченное качество открытых данных, насыщение текущих бенчмарков и отсутствие ресурсов для масштабных реальных испытаний в академических лабораториях.

\section{Обучение с подкреплением в робототехнике}
\label{sec:rl}

    Обучение с подкреплением (Reinforcement Learning, RL) -- парадигма машинного обучения, в которой агент обучается принимать решения через взаимодействие со средой, максимизируя кумулятивный сигнал вознаграждения \cite{sutton2018rl}. В контексте робототехнического манипулирования агентом выступает робот, средой -- физическое или симулированное окружение, действиями -- управляющие команды, а состояниями -- визуальные наблюдения и проприоцептивная информация.

    Среди алгоритмов RL, наиболее широко применяемых в робототехнике, выделяются PPO (Proximal Policy Optimization) \cite{schulman2017ppo} -- on-policy алгоритм с клипированным суррогатным функционалом, обеспечивающий стабильное обучение за счёт ограничения размера обновления политики, и SAC (Soft Actor-Critic) \cite{haarnoja2018sac} -- off-policy алгоритм, максимизирующий энтропийно-регуляризованный функционал, что поощряет исследование и повышает робастность к начальным условиям. TD3 (Twin Delayed DDPG) \cite{fujimoto2018td3} решает проблему переоценки Q-функции через использование двух критиков и задержанного обновления политики.

    Ключевым узким местом применения RL в робототехнике является проектирование функции вознаграждения (reward engineering). Определение информативного вознаграждения для сложных манипуляционных задач требует экспертных знаний и трудоёмкой итеративной настройки. Разреженные вознаграждения (бинарный сигнал <<задача выполнена / не выполнена>>) приводят к крайне медленному обучению, а плотные вознаграждения сложно проектировать для произвольных задач без формализации промежуточных подцелей. Эта проблема мотивирует использование визуально-языковых моделей для автоматизации проектирования вознаграждений.

\section{Интеграция RL и визуально-языковых моделей}
\label{sec:integration}

    По результатам анализа литературы выделяются три основных направления интеграции обучения с подкреплением и визуально-языковых моделей в задачах робототехнического манипулирования.

\subsection{Классификация подходов интеграции}
\label{subsec:integration-classification}

    Для систематизации работ введём классификацию по трём направлениям: (1) RL-дообучение VLA (policy fine-tuning), (2) VLM как генератор функции вознаграждения (reward generation), (3) VLM как планировщик/декомпозитор в связке с RL (planning + skills/RL). Сводная классификация приведена в таблице~\ref{tab:integration-classification}.

\begingroup
\setlength{\tabcolsep}{3pt}
\renewcommand{\arraystretch}{1.2}
\footnotesize
\begin{longtable}{>{\raggedright\arraybackslash}p{0.17\textwidth} >{\raggedright\arraybackslash}p{0.26\textwidth} >{\raggedright\arraybackslash}p{0.26\textwidth} >{\raggedright\arraybackslash}p{0.25\textwidth}}
\caption{Классификация подходов интеграции VLM/VLA и RL.}
\label{tab:integration-classification}\\
\toprule
\textbf{Аспект} & \makecell[l]{\textbf{RL-дообучение}\\\textbf{VLA}} & \makecell[l]{\textbf{VLM $\rightarrow$ награда}\\\textbf{$\rightarrow$ RL}} & \makecell[l]{\textbf{VLM-планирование}\\\textbf{+ RL/навыки}} \\
\midrule
\endfirsthead
\toprule
\textbf{Аспект} & \makecell[l]{\textbf{RL-дообучение}\\\textbf{VLA}} & \makecell[l]{\textbf{VLM $\rightarrow$ награда}\\\textbf{$\rightarrow$ RL}} & \makecell[l]{\textbf{VLM-планирование}\\\textbf{+ RL/навыки}} \\
\midrule
\endhead
\midrule
\multicolumn{4}{r}{\small Продолжение на следующей странице} \\
\endfoot
\bottomrule
\endlastfoot
\textbf{Постановка задачи} &
Улучшить готовую политику $\pi_{\text{VLA}}(a|o,l)$, повысив SR за счёт онлайнового взаимодействия (часто через residual-обёртку) &
Снять зависимость от ручной награды, сгенерировать $R_{\text{VLM}}$ из языка/видео и обучить RL-агента &
Решать длинногоризонтные задачи: VLM выбирает/синтезирует высокоуровневые подцели/планы, низкий уровень исполняется RL/навыками \\
\textbf{Модальности} &
Зрение + язык $\rightarrow$ действия; для RL-части обычно состояния/фичи среды + (иногда) визуальные наблюдения &
Зрение (кадры/рендеры) + язык (инструкция/описание цели); опционально состояние симулятора для кодогенерации &
Зрение + язык; дополнительно список навыков/примитивов, карты сцены, API среды \\
\textbf{Роль RL} &
Основной механизм улучшения политики (онлайн дообучение, residual RL, предпочтения/стадии) &
Основной механизм обучения политики при автоматически получаемой награде (код/оценка/предпочтения) &
Компонент низкоуровневого управления (навыки/контроллеры) и/или оценка выполнимости действий (affordance/value) \\
\textbf{Тип данных} &
Демонстрации для начального SFT/BC + онлайновые траектории (симуляция и/или реальность); иногда предпочтения по стадиям &
Описания задач + (часто) доступ к коду среды/состоянию + траектории RL; данные для обучения reward-модели (если есть) &
Данные для VLM (общее предобучение) + навыковые датасеты/трейсы; иногда интерактивные данные для уточнения планов \\
\textbf{Типовые метрики} &
SR/episode success, recovery success, обобщение на новые объекты/инструкции, робастность; иногда cost/latency &
SR и reward, sample efficiency (шаги до SR$_\text{target}$), стабильность по сидам; согласованность награды с успехом &
SR на многошаговых заданиях, длина плана/число подзадач, доля выполнимых действий, zero-shot SR в реальности \\
\textbf{Ограничения} &
Высокая вычислительная стоимость (крупные VLA), риски нестабильности RL-дообучения, sim-to-real разрыв, требования к безопасности &
Ошибки/галлюцинации VLM, чувствительность к API/состоянию среды, хрупкость сгенерированного кода, вычислительная цена многократных прогонов RL &
Ограниченная наблюдаемость и заземление, накопление ошибок плана, требования к библиотеке навыков, сложность оценки выполнимости \\
\textbf{Примеры работ} &
PLD \cite{pld2026}, STARE-VLA \cite{starevla2025}, VLA-RL \cite{vlarl2025}, SimpleVLA-RL \cite{simplevlarl2025}, RIPT-VLA \cite{riptvla2025}, ConRFT \cite{conrft2025} &
EUREKA \cite{ma2024eureka}, Text2Reward \cite{xie2024text2reward}, RL-VLM-F \cite{wang2024rlvlmf} &
SayCan \cite{ahn2022saycan}, Code as Policies \cite{liang2023code}, VoxPoser \cite{huang2023voxposer}, Embodied-R1 \cite{embodiedr1_2026} \\
\end{longtable}
\endgroup

\subsection{RL для дообучения VLA-моделей}
\label{subsec:rl-finetune}

    Первое направление -- применение RL для дообучения уже предобученных VLA-моделей с целью повышения их success rate. Основная идея заключается в том, что VLA-модель, обученная через имитационное обучение на демонстрациях, может быть улучшена через взаимодействие со средой при наличии сигнала вознаграждения.

    В обобщённом виде данная линия работ характеризуется следующими особенностями.
    \begin{enumerate}
        \item \textbf{Постановка задачи}: максимизация success rate путём адаптации политики, полученной через SFT/BC, к распределению состояний тестовой среды (в том числе к редким ошибочным состояниям и recovery-поведениям).
        \item \textbf{Модальности}: входы VLA обычно включают изображение/видео и текстовую инструкцию; RL-компонент может использовать как те же наблюдения, так и компактные признаки/состояния среды.
        \item \textbf{Роль RL}: дообучение (онлайн или гибридное) с целью исправления систематических ошибок и повышения устойчивости; часто применяется residual-обёртка, уменьшающая риск деградации базовой политики.
        \item \textbf{Тип данных}: демонстрации (для обучения базовой VLA) + траектории взаимодействия, собранные текущей политикой и/или residual-агентом; иногда дополнительные предпочтения или стадийные разметки.
        \item \textbf{Типовые метрики}: SR, средняя длина успешных эпизодов, показатели recovery, устойчивость по сидам; при наличии реального робота -- SR в реальности и sim-to-real разрыв.
        \item \textbf{Ограничения}: высокая вычислительная стоимость и требования к памяти (крупные VLA), потенциальная нестабильность RL-обновлений, а также сложности безопасного онлайнового обучения в реальном мире.
    \end{enumerate}

    PLD (Probe, Learn, Distill) \cite{pld2026}, представленный на ICLR 2026, реализует трёхстадийный подход. На первой стадии VLA-модель замораживается, и поверх неё обучается лёгкий residual actor через off-policy RL. Этот специалист компенсирует ошибки базовой политики в тех состояниях, где она чаще всего терпит неудачу. На второй стадии разворачивается гибридная схема сбора данных, смещающая распределение в сторону состояний, часто посещаемых базовой политикой, но при этом захватывающая recovery-поведение. На третьей стадии собранные траектории дистиллируются обратно в VLA-модель через стандартное supervised fine-tuning, поддерживая как flow-matching, так и авторегрессионные модели. Результат -- 99\% success rate на LIBERO, улучшение более чем на 50\% на SIMPLER и 100\% на реальных роботах Franka и YAM.

    STARE-VLA \cite{starevla2025} предлагает декомпозицию длинных робототехнических траекторий на семантически значимые стадии (Reach $\to$ Grasp $\to$ Transport $\to$ Place) с назначением вознаграждений на уровне каждой стадии, а не всей траектории. На основе этой идеи разработаны STA-TPO для офлайнового обучения на предпочтениях и STA-PPO для онлайнового взаимодействия со средой. Достигнуты 98\% на SIMPLER и 96,4\% на ManiSkill3.

    В 2025 году это направление развивалось особенно быстро, и появившиеся работы удобно сгруппировать по тому, как именно в них организована оптимизация политики.

    Несколько коллективов пошли по пути полного дообучения VLA на уровне токенов действий. VLA-RL \cite{vlarl2025} рассматривает траекторию манипулирования как многоходовой мультимодальный диалог и обучает авторегрессионную модель онлайновым RL; чтобы справиться с разреженностью награды, авторы добавляют отдельную reward-модель на базе VLM, обученную на псевдо-метках. На 40 задачах LIBERO поверх OpenVLA-7B это даёт прирост около 4,5~п.п. и выводит модель на уровень $\pi_0$-FAST. Близкая по духу SimpleVLA-RL \cite{simplevlarl2025} переносит на VLA простую схему из DeepSeek-R1: траектории присваивается бинарная награда за успех, после чего политика обновляется методом GRPO. Авторы сообщают о росте среднего SR на LIBERO с 91 до 99\% (на наборе Long -- с 86,5 до 98,5\%) и отмечают, что модель иногда находит стратегии, не встречавшиеся в демонстрациях.

    Другая группа методов стремится упростить или стабилизировать саму оптимизацию. RIPT-VLA \cite{riptvla2025} вообще отказывается от обучаемого критика и работает с разреженной бинарной наградой по устойчивой схеме LOOP; этого оказывается достаточно, чтобы поднять OpenVLA-OFT с 96,7 до 97,5\%, а в предельном случае одной демонстрации довести почти неработающую модель до 97\% SR. ConRFT \cite{conrft2025} сочетает офлайновое Q-обучение с онлайновым дообучением consistency-политики и участием человека, а iRe-VLA \cite{irevla2024} чередует обновления PPO при замороженном backbone с этапами supervised-дистилляции, повышая стабильность ценой необходимости в критике и shaped-награде.

    Ещё одна линия работ делает ставку на инфраструктуру обучения. VLA-RFT \cite{vlarft2025} обучает мировую модель и использует её как управляемый симулятор, получая плотную награду и обходясь менее чем 400 шагами дообучения, а RLinf-VLA \cite{rlinfvla2025} предлагает единый фреймворк сразу для нескольких архитектур, алгоритмов и симуляторов. Наконец, обстоятельное эмпирическое исследование \cite{rlvla_study2025} прямо сравнивает алгоритмы и приходит к выводу, что для VLA метод PPO в целом надёжнее, чем перенесённые из языковых моделей GRPO и DPO.

    На этом фоне выбранный в работе подход отличается тем, что базовая VLA остаётся замороженной, а обучается лишь лёгкий SAC-агент, добавляющий небольшую поправку к её действиям. Такая residual-схема ближе всего к PLD~\cite{pld2026}, но обходится без финальной дистилляции. У неё несколько практических достоинств: она не разрушает уже накопленные моделью навыки (что хорошо видно по провалу варианта с полной разморозкой), требует на порядок меньше обучаемых параметров и умеренного бюджета, и позволяет в одинаковых условиях прогнать сразу девять моделей, тогда как в большинстве упомянутых статей сравнение ограничено одной-двумя. Выбор SAC здесь тоже не случаен: поправка к действию непрерывна, и off-policy алгоритм с энтропийной регуляризацией ложится на эту задачу естественнее, чем токен-уровневые on-policy методы, удобные для авторегрессионной генерации.

    Тем не менее единого, общепринятого рецепта RL-дообучения VLA пока не сложилось, и эта проблема остаётся открытой \cite{reuss2025blog}.

\subsection{VLM как генератор функций вознаграждения для RL}
\label{subsec:vlm-reward}

    Второе направление решает проблему reward engineering за счёт автоматической генерации функций вознаграждения с помощью визуально-языковых моделей. Выделяются два основных подхода: генерация вознаграждений в виде исполняемого кода и генерация через парные сравнения наблюдений.

    Для данной линии характерны следующие общие свойства.
    \begin{enumerate}
        \item \textbf{Постановка задачи}: заменить ручное проектирование $R(s,a)$ на автоматическую процедуру, использующую описание цели на естественном языке и (при наличии) структуру среды; цель -- повысить SR и/или ускорить обучение за счёт более информативной награды.
        \item \textbf{Модальности}: язык (описание задачи, ограничения, желаемое поведение), изображения/видео (наблюдения до/после действия); в кодогенерации часто используется состояние среды и её API.
        \item \textbf{Роль RL}: обучение низкоуровневой политики в среде с наградой, полученной от VLM; в ряде работ RL выступает также как ``внутренний оценщик'' качества награды (через прогон обучения).
        \item \textbf{Тип данных}: либо (а) доступ к симулятору и его коду/параметрам для генерации награды как кода, либо (б) пары/трой\-ки наблюдений для предпочтений и последующего обучения reward-модели; дополнительно -- траектории, собранные RL-агентом.
        \item \textbf{Типовые метрики}: SR на бенчмарке (часто основная), sample efficiency (шаги/эпизоды до порога SR), дисперсия по сидам; для награды как артефакта -- доля корректно исполняемых функций и корреляция reward с успехом.
        \item \textbf{Ограничения}: хрупкость сгенерированного кода и зависимость от API, галлюцинации VLM, вычислительная стоимость итеративной оптимизации награды (много запусков RL), ограниченная переносимость на реальность без адаптации сенсоров/перцепции.
    \end{enumerate}

    EUREKA \cite{ma2024eureka} (ICLR 2024) использует GPT-4 для генерации кода функции вознаграждения на основе описания задачи и исходного кода среды. Алгоритм выполняет эволюционную оптимизацию: генерирует множество кандидатов, оценивает их через GPU-ускоренное обучение RL-агентов и итеративно улучшает через рефлексию. На 29 окружениях с 10 различными морфологиями роботов EUREKA превосходит экспертные вознаграждения на 83\% задач со средним улучшением в 52\%. Впервые был продемонстрирован pen-spinning на антропоморфной руке Shadow Hand.

    Text2Reward \cite{xie2024text2reward} (ICLR 2024 Spotlight) генерирует плотные reward-функции в виде исполняемого Python-кода, опираясь на компактное представление среды и текстовое описание цели. В отличие от EUREKA, Text2Reward работает с компактным представлением среды и поддерживает итеративное уточнение через обратную связь от пользователя. На 13 из 17 задач ManiSkill2 и MetaWorld политики, обученные с награ\-дами Text2Reward, достигают success rate, сравнимого или превосходящего экспертные. Осуществлён успешный sim-to-real перенос на реальный робот.

    RL-VLM-F \cite{wang2024rlvlmf} (ICML 2024) предлагает альтернативный подход: вместо генерации кода награды VLM используется для попарного сравнения наблюдений (<<какое из двух состояний ближе к цели?>>). Из предпочтений обучается непрерывная функция вознаграждения. Парные сравнения более устойчивы к галлюцинациям VLM, чем абсолютные числовые оценки. Метод продемонстрирован на задачах с твёрдыми, шарнирными и деформируемыми объектами.

    RoboReward \cite{roboreward2025} обучает специализированные reward-модели с 4 и 8 миллиардами параметров на робототехнических данных с counterfactual data augmentation -- генерацией calibrated negatives и near-misses. Модель 8B превосходит Gemini Robotics-ER 1.5 на задачах с коротким горизонтом.

    T2-VLM \cite{zhao2025t2vlm} (ICCV 2025) предлагает training-free подход: байесовская фильтрация используется для отслеживания VLM-генерированных подцелей и генерации темпорально-консистентных вознаграждений без дополнительного обучения.

    LaGEA \cite{lagea2025} использует schema-constrained языковые рефлексии VLM как темпорально обоснованное руководство для RL-агента. Вводятся adaptive failure-aware коэффициенты для шейпинг-наград. Результат -- улучшение на 5--17\% на MetaWorld MT10 и Fetch.

\subsection{VLM как планировщик в связке с RL}
\label{subsec:vlm-planner}

    Третье направление использует визуально-языковые модели для высокоуровневого планирования, декомпозиции задач на подцели и рассуждений, а низкоуровневое управление осуществляется RL-агентом или предобученными навыковыми политиками.

    В терминах системной архитектуры данная линия обычно включает явную иерархию.
    \begin{enumerate}
        \item \textbf{Постановка задачи}: решать композиционные и длинногоризонтные задания через декомпозицию на подзадачи и выбор примитивов/навыков в текущем контексте.
        \item \textbf{Модальности}: наблюдения (изображения/видео, иногда карты/сцена) + язык; дополнительно -- список доступных навыков/действий, их описания и ограничения.
        \item \textbf{Роль RL}: либо обучение навыков (низкий уровень), либо функция выполнимости/ценности (affordance/value) для ранжирования действий планировщика; RL может быть заменён на предобученные навыки, но в робототехнике часто используется как наиболее универсальный исполнитель.
        \item \textbf{Тип данных}: датасеты навыков/траекторий, синтетические или реальные; интерактивные эпизоды для уточнения планов; иногда -- данные VQA для замкнутого контура обратной связи.
        \item \textbf{Типовые метрики}: SR на многошаговых заданиях, показатели разложения (число шагов плана, доля валидных подзадач), устойчивость к ошибкам распознавания; при переносе -- zero-shot SR в реальности.
        \item \textbf{Ограничения}: накопление ошибок при декомпозиции, слабое заземление языка в геометрию/контакты, необходимость библиотеки навыков и интерфейса API, высокая чувствительность к перцепции и частичным наблюдениям.
    \end{enumerate}

    SayCan \cite{ahn2022saycan} объединяет языковую модель, оценивающую семантическую уместность действий, с функцией affordance, оценивающей физическую выполнимость действия в текущем состоянии. Произведение этих оценок определяет выбор следующего действия из набора предобученных навыков. Подход продемонстрирован на мобильном манипуляторе в кухонной среде.

    Code as Policies \cite{liang2023code} генерирует исполняемый Python-код политики управления на основе текстовой инструкции. Ключевой механизм -- иерархическая кодогенерация: неопределённые функции рекурсивно раскрываются в более элементарные вызовы API. Подход поддерживает spatial-geometric рассуждения и обобщение на новые инструкции.

    Inner Monologue \cite{huang2023inner} решает проблему отсутствия обратной связи в LLM-планировщиках. Модель формирует <<внутренний монолог>>, включающий описание сцены (пассивное и активное через VQA), детекцию успешности действий и рассуждения о следующем шаге. Замкнутый контур значительно улучшает выполнение многошаговых задач в реальной кухонной среде.

    VoxPoser \cite{huang2023voxposer} использует LLM для композиции 3D карт значений (value maps), определяющих области притяжения и отталкивания в трёхмерном пространстве. Карты синтезируются из текстовых описаний через визуально обоснованные пространственные рассуждения и используются для планирования траекторий.

    Embodied-R1 \cite{embodiedr1_2026} (ICLR 2026) вводит pointing как универсальное, embodiment-agnostic промежуточное представление, связывающее высокоуровневое визуально-языковое понимание с низкоуровневыми действиями. Определены четыре типа embodied pointing: REG (указание на объект), RRG (указание на место по отношению), OFG (указание на функциональную часть), VTG (последовательность точек как визуальная траектория). Модель (3B, на основе Qwen2.5-VL) обучена двухстадийным Reinforced Fine-Tuning на наборе Embodied-Points-200K. Достигнуты 87,5\% zero-shot success rate на 8 реальных задачах манипулирования роботом XArm (+62\% по сравнению с базовыми подходами).

\section{Бенчмарки}
\label{sec:benchmarks}

    Для оценки VLA-моделей и методов интеграции RL используется ряд симуляционных бенчмарков.

    LIBERO \cite{liu2023libero} содержит четыре набора задач (Goal, Spatial, Long, Object), различающихся по типу обобщения. Порог state-of-the-art составляет $>$95\% для Spatial, Goal и Object и $>$90\% для Long. На текущий момент бенчмарк практически решён: большинство современных VLA-моделей достигают 95--99\%.

    SIMPLER \cite{simpler2024} предоставляет симулированные версии реальных робототехнических задач для Bridge и Google Robot. Результаты на Bridge варьируются от 40\% до 99\%, что делает кросс-статейные сравнения затруднительными. Для Google Robot state-of-the-art составляет 70--80\%.

    MetaWorld \cite{yu2020metaworld} включает 50 задач манипулирования с оценкой в режимах MT10 (10 задач), MT50 (50 задач) и ML10/ML45 для мета-обучения. Широко используется для оценки методов VLM-генерации вознаграждений.

    ManiSkill2 \cite{gu2023maniskill2} содержит 20 семейств задач манипулирования с более чем 2000 моделями объектов и поддержкой визуального RL.

    Как отмечается в \cite{reuss2025blog}, насыщение ряда бенчмарков (особенно LIBERO) ограничивает их полезность для оценки реального прогресса. Результаты, близкие к потолку, не позволяют дифференцировать методы, а улучшения в 0,5--1\% на насыщенных бенчмарках слабо коррелируют со способностью к обобщению в реальном мире.

\section{Выводы по обзору}
\label{sec:review-summary}

    По результатам аналитического обзора можно сделать следующие выводы.

    Область VLA-моделей переживает период взрывного роста. За два года (2024--2026) число работ по тематике на ICLR выросло с единиц до 164 -- около восемнадцатикратного роста только за последний год; появились открытые модели (OpenVLA, $\pi_0$, FLOWER), сопоставимые с крупными закрытыми системами на стандартных бенчмарках, однако уступающие им в zero-shot обобщении.

    Интеграция RL и VLM/VLA развивается по трём направлениям, каждое из которых показывает значимые результаты: RL-дообучение VLA (до 99\% success rate), VLM-генерация вознаграждений (превосходство над экспертными на 83\% задач), VLM-планирование (87,5\% zero-shot в реальном мире). При этом ни одно из направлений не является универсально лучшим -- выбор зависит от задачи, модели и доступных ресурсов.

    Проблема reward engineering остаётся ключевым узким местом RL в робототехнике, и автоматизация через VLM представляет наиболее практически значимое направление для широкого круга задач.

    Текущие бенчмарки ограничены в способности оценить реальный прогресс, что требует осторожности при интерпретации результатов и мотивирует развитие новых методов оценки.

    Из всего сказанного и вырастает постановка настоящей работы. Поскольку ни одно из трёх направлений нельзя назвать безусловно лучшим, а выбор всякий раз упирается в наличие демонстраций, предобученной VLA и вычислительных ресурсов, основное внимание я сосредоточил на RL-дообучении VLA -- именно оно даёт сегодня лучшие результаты (98--99\% SR у PLD и STARE-VLA). Ветка с VLM-генерацией наград рассматривается как дополнительная, на случай, когда готовой VLA нет. Соответственно, в главе~\ref{ch:chap2} обоснована двухтрековая методика, а в главе~\ref{ch:chap3} серии~1--4 посвящены связке VLA+RL на LIBERO, и серия~5 -- VLM-наградам на MetaWorld.

\endinput